\scanpage{652}\mspage{40}
nous amène à la

\uline{Question 3.6} Soit $M$ une catégorie stable par \inserted{petites} $\lim$,
et soit $(\deleted{\unsure{C}}W,\deleted{\unsure{W}}C)$ une structure à cofibration dans
$M$ telle que $C$ et $TC=C\cap W$ soient Q-clos à gauche.
Soit $I$ un objet de Cat, a-t-on alors
\[
 C\langle I\rangle\cap W(I)\overset{?}{=}(TC)\langle I\rangle\;?\tag*{(3.6.1)}
\]

On a vu au \S\ précédent qu'il en est ainsi
si $(W,C)$ s'insère dans un triplet de
Quillen clos $(W,C,F=(TC)_*)$, \unsure{et quand on|et ayant en}
suppose de plus que les couples
\[
 C\langle I\rangle\longleftrightarrow TF(I),\qquad TC\langle I\rangle\longleftrightarrow F(I)\unsure{=}
\]
sont factorisants, auquel cas $(W(I),C\langle I\rangle,F(I))$
est une structure de Quillen close dans $M(I)$.
D'ailleurs, $(W,C,F)$ est une structure de Quillen
(néc. close) ssi les couples
$C\longleftrightarrow\underset{\inserted{\substack{\rotatebox{90}{$=$}\,\mathrm{def}\\ C_*}}}{\underbrace{TF}}$,
$TC\longleftrightarrow\underset{\inserted{\substack{\rotatebox{90}{$=$}\,\mathrm{def}\\ (TC)_*}}}{F}$
eux-mêmes sont factorisants\deleted{, et si de
plus $TF\subset W$ \unsure{(on aura alors}} et si de plus
$TF\subset W$ (ce qui est automatique si $W=W^{\unsure{\mathrm{G}|\leftrightarrow}}$,
car on aura en tous cas, par \deleted{\unsure{hypoth.}} Cof 4, $TF\subset W^{\mathrm{G}}$).

Mais peut-on s'en tirer avec moins, pour
pouvoir conclure à (3.6.1) ? Ne suffirait-il
pas de savoir que $TC\langle I\rangle\longleftrightarrow F(I)$ est
factorisant (ce qui n'a guère de chances
d'être vrai pour un $I$ $\unsure{\ne|I}$ général, si on n'a
déjà $TC\mathrel{\longleftarrow\mkern-9mu|\mkern-9mu\longrightarrow}F$, \unsure{factorisation|factorisant|factorisables} dans $M$
lui-même). C'est bien le cas :
% archive sheet number: 651

